%%
%% part3-customization/ch19-fonts.tex
%%
%% Chapter 19: Fonts — The five Persian fonts, the Latin
%% font, the math font, and the monospace font.
%%

\chapter{فونت‌ها}
\label{chap:fonts}

\section{چرا این موضوع مهم است؟}
\label{sec:fonts-why}

فونت، یکی از مهم‌ترین عوامل در خوانایی و
زیبایی کتاب است. یک فونت خوب، چشم خواننده را
خسته نمی‌کند و درک مطلب را آسان‌تر می‌سازد.

\lr{MatinBook} از پنج فونت فارسی پشتیبانی
می‌کند و همچنین فونت‌های لاتین، ریاضی و
مونواسپیس را به‌طور دقیق تنظیم می‌کند. در
این فصل، با این فونت‌ها آشنا می‌شوید و
یاد می‌گیرید که چگونه آن‌ها را شخصی‌سازی
کنید.

\mnote{\lr{MatinBook} از \lr{fontspec} و \lr{unicode-math} برای مدیریت فونت‌ها استفاده می‌کند که از فونت‌های \lr{OpenType} پشتیبانی می‌کنند.}

\section{فونت‌های فارسی}
\label{sec:fonts-persian}

\lr{MatinBook} پنج فونت فارسی را پشتیبانی
می‌کند که با آپشن کلاس قابل انتخاب هستند.
جدول \cref{tab:persian-fonts} این فونت‌ها
را خلاصه می‌کند.

\begin{matintable}{فونت‌های فارسی \lr{MatinBook}}{tab:persian-fonts}
\begin{tblr}{
    width = \textwidth,
    colspec = {l X[l] X[c] X[c]},
    hlines, vlines,
    colsep = 6pt,
    rowsep = 4pt,
    row{1} = {font=\bfseries},
}
آپشن & فونت & مقیاس & پیش‌فرض \\
\lr{niloofar} & \lr{XB Niloofar} & \pnum{۱.۲} & بله \\
\lr{vazirmatn} & \lr{Vazirmatn} & \pnum{۱.۲} & خیر \\
\lr{sahel} & \lr{Sahel} & \pnum{۱.۲} & خیر \\
\lr{irlotus} & \lr{IR Lotus} & \pnum{۱.۲} & خیر \\
\lr{bnazanin} & \lr{B Nazanin} & \pnum{۱.۰} & خیر \\
\end{tblr}
\end{matintable}

\subsection{فونت \lr{XB Niloofar}}
\label{sec:fonts-niloofar}

\lr{XB Niloofar} فونت پیش‌فرض \lr{MatinBook}
است. این فونت به‌دلیل خوانایی بالا در متن‌های
فنی و پوشش کامل حروف فارسی انتخاب شده است.

\begin{latin}
\begin{minted}{latex}
\documentclass[niloofar]{matinbook}
\end{minted}
\end{latin}

مقیاس این فونت \pnum{۱.۲} است، یعنی ۲۰٪
بزرگ‌تر از فونت پایه‌ی سند (\pnum{۱۰pt}).

\mnote{مقیاس \pnum{۱.۲} برای فونت فارسی، استاندارد نشر فنی ایران است. این تنظیم باعث می‌شود فونت فارسی با فونت لاتین تعادل بصری داشته باشد.}

\subsection{فونت \lr{Vazirmatn}}
\label{sec:fonts-vazirmatn}

\lr{Vazirmatn} یک فونت فارسی مدرن با طراحی
هندسی است که خوانایی بالایی در نمایش دیجیتال
دارد:

\begin{latin}
\begin{minted}{latex}
\documentclass[vazirmatn]{matinbook}
\end{minted}
\end{latin}

\mnote{فونت \lr{Vazirmatn} برای کتاب‌های دیجیتال و وب مناسب‌تر است، ولی در چاپ هم کیفیت خوبی دارد.}

\subsection{فونت \lr{Sahel}}
\label{sec:fonts-sahel}

\lr{Sahel} یک فونت فارسی ساده و خوانا است
که برای متن‌های طولانی مناسب است:

\begin{latin}
\begin{minted}{latex}
\documentclass[sahel]{matinbook}
\end{minted}
\end{latin}

\subsection{فونت \lr{IR Lotus}}
\label{sec:fonts-irlotus}

\lr{IR Lotus} یک فونت فارسی سنتی است که
حس کلاسیک به کتاب می‌دهد:

\begin{latin}
\begin{minted}{latex}
\documentclass[irlotus]{matinbook}
\end{minted}
\end{latin}

\subsection{فونت \lr{B Nazanin}}
\label{sec:fonts-bnazanin}

\lr{B Nazanin} یک فونت تجاری است که پوشش
حروف آن کامل نیست و برای کتاب‌های فنی
توصیه نمی‌شود:

\begin{latin}
\begin{minted}{latex}
\documentclass[bnazanin]{matinbook}
\end{minted}
\end{latin}

\mnote{فونت \lr{B Nazanin} تجاری است و ممکن است روی همه‌ی سیستم‌ها نصب نباشد. برای کتاب‌های فنی، \lr{XB Niloofar} یا \lr{Vazirmatn} توصیه می‌شود.}

\section{فونت لاتین}
\label{sec:fonts-latin}

فونت لاتین پیش‌فرض \lr{MatinBook}، \lr{Times
New Roman} است. اگر این فونت روی سیستم
نصب نباشد، \lr{TeX Gyre Termes} به‌عنوان
جایگزین استفاده می‌شود.

\begin{latin}
\begin{minted}{latex}
% |\pc{در فایل}|\lr{matinbook.cls}:
\IfFontExistsTF{Times New Roman}{
    \setlatintextfont{Times New Roman}
}{
    \setlatintextfont{TeX Gyre Termes}
    \PackageWarning{matinbook}{Times New Roman not found}
}
\end{minted}
\end{latin}

\mnote{فونت \lr{TeX Gyre Termes} یک فونت رایگان و متن‌باز است که ظاهری مشابه \lr{Times New Roman} دارد.}

\section{فونت ریاضی}
\label{sec:fonts-math}

فونت ریاضی پیش‌فرض \lr{MatinBook}، \lr{XITS
Math} است. این فونت، پوشش گسترده‌ای از
نمادهای ریاضی دارد و با \lr{unicode-math}
سازگار است.

\begin{latin}
\begin{minted}{latex}
% |\pc{در فایل}|\lr{matinbook.cls}:
\IfFontExistsTF{XITS Math}{
    \setmathfont{XITS Math}
}{
    \IfFontExistsTF{XITSMath-Regular.otf}{
        \setmathfont{XITSMath-Regular.otf}
    }{
        \IfFontExistsTF{Latin Modern Math}{
            \setmathfont{Latin Modern Math}
        }{
            \PackageWarning{matinbook}{No math font found}
        }
    }
}
\end{minted}
\end{latin}

\mnote{فونت \lr{XITS Math} از \lr{OpenType} پشتیبانی می‌کند و نمادهای ریاضی را با کیفیت بالا نمایش می‌دهد.}

\section{فونت مونواسپیس}
\label{sec:fonts-monospace}

فونت مونواسپیس پیش‌فرض \lr{MatinBook}،
\lr{DejaVu Sans Mono} است. این فونت برای
نمایش کد، اعداد و متن‌های فنی مناسب است.

\begin{latin}
\begin{minted}{latex}
% |\pc{در فایل}|\lr{matinbook.cls}:
\IfFontExistsTF{DejaVu Sans Mono}{
    \setmonofont{DejaVu Sans Mono}[Scale=0.9]
}{
    \IfFontExistsTF{Courier New}{
        \setmonofont{Courier New}[Scale=0.9]
    }{
        \PackageWarning{matinbook}{No monospace font found}
    }
}
\end{minted}
\end{latin}

\mnote{فونت \lr{DejaVu Sans Mono} روی اکثر توزیع‌های \lr{Linux} به‌طور پیش‌فرض نصب است.}

\section{فونت کد}
\label{sec:fonts-code}

فونت کد پیش‌فرض \lr{MatinBook}، \lr{JetBrains
Mono} است. این فونت، لیگچرهای مخصوص
برنامه‌نویسی دارد که خوانایی کد را افزایش
می‌دهند.

\begin{latin}
\begin{minted}{latex}
% |\pc{در فایل}|\lr{matinbook.cls}:
\IfFontExistsTF{JetBrains Mono}{
    \newfontfamily\codefont{JetBrains Mono}[Scale=0.9,Ligatures=TeX]
}{
    \IfFontExistsTF{DejaVu Sans Mono}{
        \newfontfamily\codefont{DejaVu Sans Mono}[Scale=0.9,Ligatures=TeX]
    }{
        \newfontfamily\codefont{Courier New}[Scale=0.9]
    }
}
\end{minted}
\end{latin}

\mnote{اگر \lr{JetBrains Mono} نصب نباشد، \lr{DejaVu Sans Mono} به‌عنوان جایگزین استفاده می‌شود.}

\section{فونت فارسی در کد}
\label{sec:fonts-persian-code}

برای نمایش متن فارسی داخل کد، \lr{MatinBook}
از فونت \lr{DejaVu Sans Mono} استفاده می‌کند
که از حروف فارسی پشتیبانی می‌کند.

\begin{latin}
\begin{minted}{latex}
% |\pc{در فایل}|\lr{mb-code.sty}:
\defpersianfont\persiancodingfont{DejaVu Sans Mono}[Script=Arabic,Scale=0.9]
\newcommand{\pc}[1]{\rl{\persiancodingfont\rl{#1}}}
\end{minted}
\end{latin}

\mnote{دستور \lr{\textbackslash pc} برای متن فارسی غیرکامنت و \lr{\textbackslash pcc} برای کامنت فارسی در پایتون استفاده می‌شود.}

\section{نصب فونت‌ها}
\label{sec:fonts-installation}

\subsection{نصب \lr{XB Niloofar}}
\label{sec:fonts-install-niloofar}

\begin{latin}
\begin{minted}{bash}
mkdir -p ~/.local/share/fonts
cp fonts/niloofar/*.ttf ~/.local/share/fonts/
fc-cache -fv
\end{minted}
\end{latin}

\subsection{نصب \lr{XITS Math}}
\label{sec:fonts-install-xits}

\begin{latin}
\begin{minted}{bash}
# Debian/Ubuntu
sudo apt install fonts-xits

# macOS
brew install --cask font-xits
\end{minted}
\end{latin}

\subsection{نصب \lr{DejaVu Sans Mono}}
\label{sec:fonts-install-dejavu}

\begin{latin}
\begin{minted}{bash}
# Debian/Ubuntu
sudo apt install fonts-dejavu
\end{minted}
\end{latin}

\mnote{برای نصب فونت‌ها روی ویندوز، روی فایل \lr{.ttf} راست‌کلیک کرده و \lr{Install} بزنید.}

\section{تغییر فونت پیش‌فرض}
\label{sec:fonts-change-default}

اگر می‌خواهید فونت پیش‌فرض را برای همیشه
تغییر دهید، فایل \file{matinbook.cls} را
ویرایش کنید:

\begin{latin}
\begin{minted}{latex}
% |\pc{مقدار پیش‌فرض}|\lr{XBNiloofar}|%
\def\matin@fontname{XBNiloofar}

% |\pc{تغییر به}|\lr{Vazirmatn}|%
\def\matin@fontname{Vazirmatn}
\end{minted}
\end{latin}

\mnote{با این تغییر، فونت پیش‌فرض کلاس \lr{Vazirmatn} می‌شود و نیازی به آپشن در هر سند نیست.}

\section{خطاهای رایج}
\label{sec:fonts-mistakes}

\subsection{فونت پیدا نشد}
\label{sec:fonts-mistake-notfound}

\textbf{مشکل:} اگر فونت روی سیستم نصب نباشد،
خطای \lr{Font X not found} می‌دهد.

\textbf{راه‌حل:} فونت را نصب کنید و
\lr{fc-cache -fv} را اجرا کنید.

\subsection{فونت تجاری}
\label{sec:fonts-mistake-commercial}

\textbf{مشکل:} فونت \lr{B Nazanin} تجاری
است و ممکن است روی همه‌ی سیستم‌ها نصب
نباشد.

\textbf{راه‌حل:} از فونت‌های رایگان مانند
\lr{XB Niloofar} یا \lr{Vazirmatn}
استفاده کنید.

\subsection{فونت فارسی در کد}
\label{sec:fonts-mistake-persian-code}

\textbf{مشکل:} اگر برای متن فارسی در کد از
فونت لاتین استفاده کنید، نتیجه مربع‌های
خالی می‌شود.

\textbf{راه‌حل:} از دستورات \cmd{pc} و
\cmd{pcc} استفاده کنید که فونت فارسی
مناسب را فعال می‌کنند.

\section{تمرین‌ها}
\label{sec:fonts-exercises}

\begin{exercise}
    یک کتاب با فونت \lr{Vazirmatn} بسازید و
    آن را با \lr{XB Niloofar} مقایسه کنید.
    \label{exc:fonts-vazirmatn}
\end{exercise}

\begin{exercise}
    فونت پیش‌فرض کلاس را به \lr{Sahel}
    تغییر دهید.
    \label{exc:fonts-default}
\end{exercise}

\begin{exercise}
    یک نمونه کد با کامنت فارسی بنویسید و
    فونت \lr{DejaVu Sans Mono} را ببینید.
    \label{exc:fonts-code}
\end{exercise}

\begin{exercise}
    نصب بودن فونت‌های اصلی را با \lr{fc-list}
    بررسی کنید.
    \label{exc:fonts-check}
\end{exercise}

\section{خلاصه}
\label{sec:fonts-summary}

در این فصل، با فونت‌های \lr{MatinBook}
آشنا شدیم:

\begin{itemize}
    \item \lr{MatinBook} پنج فونت فارسی
    را پشتیبانی می‌کند: \lr{niloofar}،
    \lr{vazirmatn}، \lr{sahel}، \lr{irlotus}،
    \lr{bnazanin}.

    \item فونت پیش‌فرض، \lr{XB Niloofar} با
    مقیاس \pnum{۱.۲} است.

    \item فونت لاتین، \lr{Times New Roman}
    (با جایگزین \lr{TeX Gyre Termes}).

    \item فونت ریاضی، \lr{XITS Math}.

    \item فونت مونواسپیس، \lr{DejaVu Sans Mono}.

    \item فونت کد، \lr{JetBrains Mono} با
    لیگچرهای برنامه‌نویسی.

    \item متن فارسی در کد با \lr{DejaVu Sans Mono}
    و دستورات \cmd{pc} و \cmd{pcc} نمایش داده
    می‌شود.

    \item خطاهای رایج شامل فونت پیدا نشد،
    فونت تجاری، و فونت فارسی در کد است.
\end{itemize}

در فصل بعدی (\cref{chap:layout})، با
صفحه‌آرایی در \lr{MatinBook} آشنا می‌شوید
و یاد می‌گیرید که چگونه چیدمان کتاب را
شخصی‌سازی کنید.